\documentclass[11pt,a4paper]{article}
\usepackage[margin=25mm]{geometry}
\usepackage{amsmath,amssymb,booktabs,enumitem}
\usepackage[T1]{fontenc}
\usepackage{lmodern}
\usepackage[colorlinks=true,linkcolor=blue,urlcolor=blue,citecolor=blue]{hyperref}
\setlength{\parindent}{0pt}
\setlength{\parskip}{6pt}
\setlist{nosep,leftmargin=*}
\newcommand{\Prob}{\mathbb{P}}
\title{Conditional waiting times:\\from a three-state CTMC to GeoSSE}
\author{A worked note accompanying the DivExtra GeoSSE demonstration}
\date{14 October 2026}
\begin{document}
\maketitle

\section{Start here: success and waiting}
We start in B at time 0. Define:
\begin{itemize}
\item \textbf{Success:} the state at time 1 is A.
\item \textbf{Waiting:} no change occurs during the first \(w\) units.
\end{itemize}
We want:
\[
S(w)=\Prob(\text{waiting}\mid\text{success}).
\]
By the definition of conditional probability,
\[
\Prob(\text{waiting}\mid\text{success})
=\frac{\Prob(\text{waiting AND success})}{\Prob(\text{success})}.
\]
\textbf{The denominator:} starting from B, the probability of ending in A
after one unit is, by definition,
\[
\Prob(\text{success})=h_B(1).
\]
\textbf{The numerator:} to both wait and succeed, two successive things must happen:
\begin{enumerate}
\item \textbf{No change for \(w\) units.} The total rate of leaving B is 2,
so this has probability \(e^{-2w}\).
\item \textbf{Then end in A after the remaining \(1-w\) units.}
Because no change occurred, we are still in B. The Markov property says the
remaining success probability is \(h_B(1-w)\).
\end{enumerate}
Multiply these successive probabilities:
\[
\Prob(\text{waiting AND success})=e^{-2w}h_B(1-w).
\]
Divide by the probability of success:
\begin{equation}
\boxed{S(w)=\frac{e^{-2w}h_B(1-w)}{h_B(1)}}.
\label{eq:key}
\end{equation}
\textbf{In words:}
\begin{quote}
Among all histories that end in A, what fraction spent their first \(w\)
units without changing?
\end{quote}
That fraction is exactly \(S(w)\). We are \textbf{not assuming waiting and
success are independent}---the shorter remaining time is precisely why the
numerator uses \(h_B(1-w)\), rather than \(h_B(1)\).

\newpage
\section{Numbers and the next-event draw}
Now make the underlying three-state continuous-time Markov chain (CTMC)
explicit. With rows and columns ordered A, B, C, its rate matrix is
\[
Q=\begin{pmatrix}-2&1&1\\1&-2&1\\1&1&-2\end{pmatrix}.
\]
The total rate of leaving any state is 2. More generally, write
\[
h_i(t)=\Prob\{\text{end in A after }t\text{ more units}\mid\text{currently in }i\}.
\]
Here \(t\) is \emph{time remaining}, not elapsed time. This probability includes
any number of intervening changes, including departures from A followed by returns;
it is not the probability of merely visiting A before the deadline.
For compact notation below, let \(C\) denote success and \(N_w\) denote waiting.
Every probability also conditions on starting in B. Waiting excludes departures
and returns to B during the interval.

For this symmetric example, \(P(t)=\exp(Qt)\), giving
\[
h_A(t)=\frac{1+2e^{-3t}}{3},\qquad
h_B(t)=h_C(t)=\frac{1-e^{-3t}}{3}.
\]
For a half-unit wait,
\[
\Prob(N_{0.5})=e^{-1}\simeq0.368,\qquad
\Prob(N_{0.5}\mid C)
=e^{-1}\frac{1-e^{-1.5}}{1-e^{-3}}\simeq0.301.
\]
Among 10,000 ordinary histories, the expected number ending in A is approximately
3,167. The expected number both waiting half a unit and ending in A is approximately
953. Thus the fraction waiting among successful histories is about
\(953/3167=0.301\). These are expected counts, not simulation results.

\subsection{From a no-event probability to a waiting time}
Let \(W\) be the next-change waiting time. Then \(S(w)=\Prob(W>w\mid C)\)
for intervals before the endpoint. Define the \emph{cumulative hazard}
\[
H(w)=-\log S(w)=2w+\log h_B(1)-\log h_B(1-w).
\]
Draw \(Z\sim\operatorname{Exp}(1)\), and find the first \(w\) for which
\(H(w)=Z\). This works because
\[
\Prob\{Z>H(w)\}=e^{-H(w)}=S(w).
\]
The exponential draw is a dimensionless threshold; it is \emph{not} a waiting
time drawn at the original event rate. Equivalently, draw \(U\) uniformly on
\((0,1)\) and solve \(S(w)=U\).

Starting in B and ending in A, \(h_B(0)=0\), so the probability of no change
throughout the remaining interval is zero. Starting in A, that probability can
be positive. If the threshold exceeds the cumulative hazard to the endpoint,
record no further event. Do not force an event in this case.

\subsection{Choose the destination, then repeat}
If a change occurs with \(u\) time units remaining, choose \(j\ne i\) with
weights \(q_{ij}h_j(u)\). From B with \(u=0.5\), this gives approximately
65\% for A and 35\% for C. Recompute the waiting-time distribution from the
new state and remaining time. Reaching A does not stop the process: it can
leave and return before the endpoint.

Both timing and destination must be conditioned. Weighting destinations but
retaining ordinary exponential waiting times is generally incorrect. The
conditional instantaneous rates are
\[
q^*_{ij}(u)=q_{ij}\frac{h_j(u)}{h_i(u)},\qquad j\ne i,
\]
where \(h_i(u)>0\). This is a Doob \(h\)-transform. Even a constant original
\(Q\) gives time-dependent conditional rates. Endpoint-conditioned CTMC sampling,
including direct, rejection and uniformization methods, is reviewed by
Hobolth and Stone~\cite{hobolth}.

\newpage
\section{The same reasoning in GeoSSE}
GeoSSE has three geographic states: A, B and AB. In addition to range changes,
lineages speciate and go extinct. We must distinguish two required outcomes:
\begin{description}
\item[\(D_i(t)\):] The likelihood of the observed descendants of a particular
backbone branch, given a lineage in state \(i\) at age \(t\).
For a terminal branch ending in A under complete sampling, this is the probability
of exactly one living descendant, in A, and no other living descendants.
\item[\(E_i(t)\):] The probability that a lineage in state \(i\) at age \(t\)
leaves no living descendants at the present. Descendants can change state and
speciate; the entire descendant clade must eventually die out.
\end{description}
For a present-day A tip and literal extinction conditioning,
\[
(D_A(0),D_B(0),D_{AB}(0))=(1,0,0),\qquad E_i(0)=0.
\]
The \(E\) curves depend on state, time and model parameters, not on which
observed branch is being traversed. The \(D\) curves depend on its observed
descendants: an A-ending branch and a B-ending branch generally have different
\(D\) curves. Their entries are likelihoods, not normalized state probabilities.

\subsection{Waiting along a backbone or a side lineage}
Let \(t\) be the current age, \(u<t\) a younger age, and \(r_i(v)\) the
\emph{original} total event rate in state \(i\). The same conditional-probability
argument gives
\begin{equation}
S_i(t,u)=\exp\left[-\int_u^t r_i(v)\,dv\right]\frac{h_i(u)}{h_i(t)},
\quad h=D\text{ on the backbone},\quad h=E\text{ on an extinct side lineage}.
\end{equation}
For constant rates this integral is \(r_i(t-u)\). With epochs, integrate each
epoch's rate over its overlap with the interval. The conditional cumulative hazard is
\[
H_i(t,u)=\int_u^t r_i(v)\,dv+\log h_i(t)-\log h_i(u).
\]
The R variable called \texttt{hazard} in the demonstration holds this
\emph{cumulative} hazard, not the instantaneous rate. An \(\operatorname{Exp}(1)\)
threshold is inverted to obtain the event age. Age decreases during simulation.
For a backbone branch, inversion stops at its younger node or tip. A side
clade is instead conditioned on extinction by age zero.

\subsection{Choosing events at the sampled age}
Write \(\lambda_{i\to\{j,k\}}\) for the rate of one unordered splitting
event. At a hidden backbone split, the two assignments have weights
\[
\underbrace{\lambda_{i\to\{j,k\}}D_jE_k}_{j\text{ continues; }k\text{ dies out}},
\qquad
\underbrace{\lambda_{i\to\{j,k\}}D_kE_j}_{k\text{ continues; }j\text{ dies out}}.
\]
Range-change weights are \(q_{ij}D_j\); backbone death is incompatible with
the required descendants. If \(j=k\), the combined hidden-split weight is
\(2\lambda D_jE_j\): either daughter can carry the observed descendants.

\newpage
\section{Explicitly expanding extinct species}
For an extinction-conditioned side lineage, the unnormalized event weights are
\[
\begin{array}{ll}
\text{Death of the lineage:} & \mu_i,\\[2pt]
\text{Range change }i\to j: & q_{ij}E_j(t),\\[2pt]
\text{Split }i\to\{j,k\}: & \lambda_{i\to\{j,k\}}E_j(t)E_k(t).
\end{array}
\]
Normalize over the available events after drawing the event age. After a split,
continue both daughters, each conditional on extinction. There is no factor of
two here: both daughters must die out, and the split is counted once.

For example, an A side lineage has weights
\[
A\to\varnothing:\ x_A,\qquad
A\to AB:\ d_AE_{AB}(t),\qquad
A\to A+A:\ s_AE_A(t)^2.
\]
An AB lineage does not die immediately in standard GeoSSE: local loss of A
produces B, and local loss of B produces A. It can subsequently die through
the process starting from those narrower ranges.

\section{The two-tip demonstration}
The accompanying \texttt{geosse-two-tip-walkthrough.R} uses a 5-Ma root,
one observed A tip and one observed B tip, complete sampling, and known rates
\[
(s_A,s_B,s_{AB},x_A,x_B,d_A,d_B)
=(0.21,0.16,0.09,0.13,0.11,0.07,0.05).
\]
It integrates \texttt{diversitree:::derivs.geosse} backward, combines the two
daughter likelihoods at the observed root, and checks the log likelihood
against \texttt{make.geosse}. Both give approximately \(-5.734600\), with
\texttt{condition.surv=FALSE} and the specified root prior. This disables an
additional crown-survival normalization; it does not remove conditioning of
the mapping on the observed tree and tips.

Root states have weights \(\pi_iD_i(\text{root})\). At the observed root,
daughter states are sampled jointly using the splitting rate and both daughter
likelihoods. Unlike a hidden split, both daughters are observed. Diversitree
assigns half of an unordered mixed-state splitting rate to each orientation
on the two labelled observed branches.

The deliberately selected short realization has these events:
\begin{center}
\begin{tabular}{rl}
\toprule
Age (Ma)&Event\\\midrule
5.000&Observed root: \(AB\to AB+B\)\\
4.535&Left branch: \(AB\to A\)\\
2.938&Hidden split: \(A\to A+A\)\\
1.432&Side daughter dies\\
0.000&Observed A and B tips survive\\\bottomrule
\end{tabular}
\end{center}
This is an illustration, not an estimate of the expected number of extinct
species. The script exposes each waiting-time draw and event-weight table.

\newpage
\section{What the construction does---and does not---claim}
\begin{enumerate}
\item \textbf{Avoiding rejection.} Ordinary simulation proposes histories
without considering their endpoint. Conditional timing and event weights
sample directly from compatible histories. The no-event formula alone is
not sufficient: the event-choice weights must also be correct.
\item \textbf{Not forcing a state or killing survivors.} The required outcome
changes the distribution throughout the history. We do not change the last
state by hand or remove living side descendants afterward.
\item \textbf{Complete versus incomplete sampling.} In this note,
\(E_i(0)=0\) means literal extinction. In an incomplete-sampling likelihood,
\(E_i(0)=1-\rho_i\) instead describes no \emph{sampled} living descendants;
unobserved survivors are possible. These are different conditional targets.
\item \textbf{No fossils.} Simulated extinct species are latent histories,
not fossil observations. Conditioning on fossils requires an appropriate
fossil observation model.
\item \textbf{Numerical accuracy.} The conditional-probability identity is
mathematical. An implementation still approximates probability curves and
inverts them numerically. It must be checked against analytical examples,
independent rejection samplers where feasible, and numerical refinement.
\item \textbf{Parameter uncertainty.} The example fixes all rates. Variation
among maps therefore describes conditional history uncertainty, not uncertainty
in the estimated rates or model choice.
\end{enumerate}

\subsection{Connection to Freyman and H\"ohna}
Freyman and H\"ohna~\cite{freyman} use backward descendant likelihoods and
forward-time equations to sample state histories along observed branches over
short time intervals. The common principle is to weight forward possibilities
by compatibility with observed descendants. The event-time construction here
instead inverts conditional no-event probabilities and subsequently expands
extinction-conditioned side clades explicitly. It is not a verbatim
implementation of their forward-equation algorithm. Agreement of backbone
distributions under matched models and conditioning is a validation target,
not a validation result established by this note.

\subsection{One sentence to remember}
\begin{quote}
The conditional probability of waiting equals the ordinary probability of
waiting, multiplied by how that wait changes the probability of completing
the required history.
\end{quote}

\begin{thebibliography}{9}
\bibitem{hobolth}
Hobolth, A. and Stone, E. A. (2009).
Simulation from endpoint-conditioned, continuous-time Markov chains on a finite
state space, with applications to molecular evolution.
\emph{The Annals of Applied Statistics}, 3, 1204--1231.
\href{https://doi.org/10.1214/09-AOAS247}{doi:10.1214/09-AOAS247}.

\bibitem{freyman}
Freyman, W. A. and H\"ohna, S. (2019).
Stochastic character mapping of state-dependent diversification reveals the
tempo of evolutionary decline in self-compatible Onagraceae lineages.
\emph{Systematic Biology}, 68, 505--519.
\href{https://pubmed.ncbi.nlm.nih.gov/30476308/}{Publication record}.
\end{thebibliography}
\end{document}
